% Section A.3, the solutions to the hand-training exercises in Chapter 6, from the walk-through in
% lectures/L06/appendix/b_exercises.md.

\section{Chapter 6: Convolutional Neural Networks}
\label{sol:sec:c6}

\solution{c6:ex:1}{Feedforward}
The whole feedforward pass is worked through below, step by step.

\step{a) Padding the input image}
The kernel is $2 \times 2$, so $n = 2 / 2 = 1$: one zero on each side. The padded image $P$ is
$6 \times 6$:
\[
  P = \begin{bmatrix}
    0 & 0 & 0 & 0 & 0 & 0 \\
    0 & 1 & 1 & 1 & 1 & 0 \\
    0 & 1 & 0 & 0 & 1 & 0 \\
    0 & 1 & 0 & 0 & 1 & 0 \\
    0 & 1 & 1 & 1 & 1 & 0 \\
    0 & 0 & 0 & 0 & 0 & 0
  \end{bmatrix}
\]

\step{b) Feedforward through the conv layer}
Output cell $(i, j)$ is computed from the $2 \times 2$ window of $P$ whose top-left corner is at
$(i, j)$, for $i$ and $j$ from 0 to 3. Two cells as examples, the top-left and the bottom-right:
\begin{align*}
  \text{sum}(0,0) &= 0 \cdot 0.2 + 0 \cdot 0.4 + 0 \cdot 0.6 + 1 \cdot 0.8 + 0.5 = 1.3 \\
  \text{sum}(3,3) &= 0 \cdot 0.2 + 1 \cdot 0.4 + 1 \cdot 0.6 + 1 \cdot 0.8 + 0.5 = 2.3
\end{align*}
Once the kernel and bias have been applied across the entire image, we get the following conv
output ($4 \times 4$):
\[
  \begin{bmatrix}
    1.3 & 1.9 & 1.9 & 1.9 \\
    1.7 & 1.7 & 1.1 & 1.9 \\
    1.7 & 1.3 & 0.5 & 1.7 \\
    1.7 & 2.1 & 1.9 & 2.3
  \end{bmatrix}
\]
Every sum is positive, the smallest being 0.5, so ReLU passes each of them through unchanged: the
pre-activation sums and the output are the same matrix here.

\step{c) The max pooling layer (feedforward)}
We split the conv layer's output into four pools.

\textbf{Top-left corner:} max value 1.9, at $(0, 1)$:
\[
  \begin{bmatrix} 1.3 & 1.9 \\ 1.7 & 1.7 \end{bmatrix}
\]
\textbf{Top-right corner:} the max value 1.9 occurs in three places; we pass forward the first
instance, at $(0, 2)$:
\[
  \begin{bmatrix} 1.9 & 1.9 \\ 1.1 & 1.9 \end{bmatrix}
\]
\textbf{Bottom-left corner:} max value 2.1, at $(3, 1)$:
\[
  \begin{bmatrix} 1.7 & 1.3 \\ 1.7 & 2.1 \end{bmatrix}
\]
\textbf{Bottom-right corner:} max value 2.3, at $(3, 3)$:
\[
  \begin{bmatrix} 0.5 & 1.7 \\ 1.9 & 2.3 \end{bmatrix}
\]
The max pooling layer's output is therefore:
\[
  \begin{bmatrix} 1.9 & 1.9 \\ 2.1 & 2.3 \end{bmatrix}
\]

\step{d) Flatten (feedforward)}
The max pooling layer's output is flattened, row by row, into a vector so it can be fed to the next
layer:
\[
  [1.9, 1.9, 2.1, 2.3]
\]

\solution{c6:ex:2}{Backpropagation and optimization}
The dense layer sends back the gradients $[10, 20, 30, 40]$.

\step{a) Flatten (backpropagation)}
The gradients from the dense layer are reshaped back into a matrix, row by row, the exact inverse
of the flatten in Exercise~\ref{c6:ex:1}:
\[
  \begin{bmatrix} 10 & 20 \\ 30 & 40 \end{bmatrix}
\]

\step{b) Max pool (backpropagation)}
Each gradient goes to the position its pool's maximum came from, found in
Exercise~\ref{c6:ex:1}, and the remaining positions get 0.

\textbf{Top-left corner:} gradient 10 is propagated to the max value 1.9 (top right):
\[
  \begin{bmatrix} 0 & 10 \\ 0 & 0 \end{bmatrix}
\]
\textbf{Top-right corner:} gradient 20 is propagated to the first occurrence of the max value 1.9
(top left):
\[
  \begin{bmatrix} 20 & 0 \\ 0 & 0 \end{bmatrix}
\]
\textbf{Bottom-left corner:} gradient 30 is propagated to the max value 2.1 (bottom right):
\[
  \begin{bmatrix} 0 & 0 \\ 0 & 30 \end{bmatrix}
\]
\textbf{Bottom-right corner:} gradient 40 is propagated to the max value 2.3 (bottom right):
\[
  \begin{bmatrix} 0 & 0 \\ 0 & 40 \end{bmatrix}
\]
Assembling all the pools' gradients into a single matrix gives us the gradient matrix that arrives
at the conv layer's output:
\[
  \text{grad} = \begin{bmatrix}
    0 & 10 & 20 & 0 \\
    0 & 0 & 0 & 0 \\
    0 & 0 & 0 & 0 \\
    0 & 30 & 0 & 40
  \end{bmatrix}
\]

\step{c) Conv layer (backpropagation)}
Every pre-activation sum from Exercise~\ref{c6:ex:1} is positive and the activation is ReLU, so
$\sigma'(s) = 1$ at every position, and multiplying each incoming gradient by it changes nothing.

\textbf{Bias gradient:} the sum of all values in the gradient matrix:
\[
  \text{biasGradient} = 0 + 10 + 20 + 0 + 0 + 0 + 0 + 0 + 0 + 0 + 0 + 0 + 0 + 30 + 0 + 40 = 100
\]

\textbf{Kernel gradients:} only four gradients are non-zero, at $(0, 1)$, $(0, 2)$, $(3, 1)$ and
$(3, 3)$, so each kernel element gets four terms, one per non-zero gradient, reading the padded
input $P$ from part (a) at that kernel element's offset:
\begin{align*}
  k_{00} &= P[0][1] \cdot 10 + P[0][2] \cdot 20 + P[3][1] \cdot 30 + P[3][3] \cdot 40 \\
         &= 0 \cdot 10 + 0 \cdot 20 + 1 \cdot 30 + 0 \cdot 40 = 30 \\
  k_{01} &= P[0][2] \cdot 10 + P[0][3] \cdot 20 + P[3][2] \cdot 30 + P[3][4] \cdot 40 \\
         &= 0 \cdot 10 + 0 \cdot 20 + 0 \cdot 30 + 1 \cdot 40 = 40 \\
  k_{10} &= P[1][1] \cdot 10 + P[1][2] \cdot 20 + P[4][1] \cdot 30 + P[4][3] \cdot 40 \\
         &= 1 \cdot 10 + 1 \cdot 20 + 1 \cdot 30 + 1 \cdot 40 = 100 \\
  k_{11} &= P[1][2] \cdot 10 + P[1][3] \cdot 20 + P[4][2] \cdot 30 + P[4][4] \cdot 40 \\
         &= 1 \cdot 10 + 1 \cdot 20 + 1 \cdot 30 + 1 \cdot 40 = 100
\end{align*}
The kernel gradients are therefore:
\[
  \begin{bmatrix} 30 & 40 \\ 100 & 100 \end{bmatrix}
\]

\textbf{Input gradients (padded input):} each non-zero gradient spreads the kernel's weights,
multiplied by that gradient, over the $2 \times 2$ patch of a new $6 \times 6$ matrix $dX$ it came
from. Only four gradients are non-zero, so only their contributions need computing:
\begin{itemize}
  \item From \code{grad[0,1] = 10}: \code{dX(0,1) += 2}, \code{dX(0,2) += 4}, \code{dX(1,1) += 6},
    \code{dX(1,2) += 8}.
  \item From \code{grad[0,2] = 20}: \code{dX(0,2) += 4}, \code{dX(0,3) += 8}, \code{dX(1,2) += 12},
    \code{dX(1,3) += 16}.
  \item From \code{grad[3,1] = 30}: \code{dX(3,1) += 6}, \code{dX(3,2) += 12}, \code{dX(4,1) += 18},
    \code{dX(4,2) += 24}.
  \item From \code{grad[3,3] = 40}: \code{dX(3,3) += 8}, \code{dX(3,4) += 16}, \code{dX(4,3) += 24},
    \code{dX(4,4) += 32}.
\end{itemize}
After summing all contributions, we get the padded input gradient matrix:
\[
  dX = \begin{bmatrix}
    0 & 2 & 8 & 8 & 0 & 0 \\
    0 & 6 & 20 & 16 & 0 & 0 \\
    0 & 0 & 0 & 0 & 0 & 0 \\
    0 & 6 & 12 & 8 & 16 & 0 \\
    0 & 18 & 24 & 24 & 32 & 0 \\
    0 & 0 & 0 & 0 & 0 & 0
  \end{bmatrix}
\]
Removing the padding, the outermost row and column on every side, leaves a $4 \times 4$ matrix
matching the original image; this is the gradient with respect to the input, which is passed further
back through the network:
\[
  \begin{bmatrix}
    6 & 20 & 16 & 0 \\
    0 & 0 & 0 & 0 \\
    6 & 12 & 8 & 16 \\
    18 & 24 & 24 & 32
  \end{bmatrix}
\]

\step{d) Optimization}
Finally, the kernel and bias are updated using the learning rate $\LR = 0.001$:
\begin{align*}
  k_{00} &= 0.2 + 0.001 \cdot 30 = 0.23 &
  k_{01} &= 0.4 + 0.001 \cdot 40 = 0.44 \\
  k_{10} &= 0.6 + 0.001 \cdot 100 = 0.70 &
  k_{11} &= 0.8 + 0.001 \cdot 100 = 0.90 \\
  \text{bias} &= 0.5 + 0.001 \cdot 100 = 0.6
\end{align*}
\textbf{Updated kernel:}
\[
  \begin{bmatrix} 0.23 & 0.44 \\ 0.70 & 0.90 \end{bmatrix}
\]
\textbf{Updated bias:} 0.6.

These are the values the demo in Exercise~\ref{c6:ex:3} prints in its \emph{expected} column: the
conv output from Exercise~\ref{c6:ex:1}, and the input gradients, the kernel and the bias from this
one.
